

\paragraph{Measure deferred balances from explicit recognition assumptions.}
Carrying a recorded deferred-tax asset and liability preserves the book, but a
later assessment must explain whether their amounts should change. The mixed
book records an asset of 10 and a liability of 4. Its current-tax declarations
also contain tax bases. Neither fact establishes the amount of deductions that
can be recognized, the applicable rate on recovery or settlement, or whether a
movement belongs in profit or loss. The conditional deferred-assessment method
therefore requires those decisions separately and binds them to the actual
reassessed book.

The distinction follows the historical source treatment. IAS 12 paragraphs 7--9
distinguish asset and liability tax bases; paragraphs 24 and 27--29 condition
asset recognition on available taxable profit and its relevant grouping and
timing. Paragraphs 47 and 51 connect measurement to the applicable rate and
manner of recovery or settlement. Paragraphs 58 and 61A distinguish profit-or-loss
movements from tax related to items outside profit or loss
\citep[IAS 12, paras. 7--9, 24, 27--29, 47, 51, 58 and 61A]{eu20231803}.
These retained historical provisions explain the method's boundaries. The
example's rates, bases and recovery capacity remain explicit hypothetical
premises, rather than findings about an actual company's tax treatment.

The selected model uses one declared entity, tax authority and recovery group,
and keeps gross deferred assets and liabilities separate. Every nondeferred
asset and liability account is represented. Inventory is divided between other
inventory and retained lots; lease assets and liabilities are divided by right;
each remaining tax-loss lot is included separately. A component and its aggregate
cannot both contribute the same difference. Item bases must reconcile with the
complete prior declarations, while recovered-lot bases must also agree with
component memory. A finding about the absence of unrepresented tax-only items
and credits makes this finite scope explicit.

Let $C_i$ be the positive magnitude of an item's carrying amount, $B_i$ its
declared tax base and $r_i$ its declared nonnegative rate. An asset has signed
taxable difference $C_i-B_i$; a liability has $B_i-C_i$. A positive difference
is taxable and a negative one is deductible. For a loss lot, the supplied model
uses the negative of its remaining quantity, without consuming that quantity.
Write $d_i$ for the nonnegative deductible difference and $t_i$ for the
nonnegative taxable difference. The caller supplies a recognized deduction
$u_i$ and a common recovery capacity $K$, all in the same tax-base currency units:
\[
 0\leq u_i\leq d_i,\qquad \sum_i u_i\leq K,\qquad
 A_1=\sum_i r_i u_i,\qquad L_1=\sum_i r_i t_i.
\]
The sums include deductible temporary differences and loss lots together, so
one capacity cannot be reused for each item. The unrecognized amount $d_i-u_i$
remains visible. A deferred-tax liability does not automatically enlarge $K$:
the timing, entity and income restrictions needed to support recovery must be
supplied explicitly. The method does not produce that support by forecasting
future taxable profit.

Before changing balances, the request allocates the recorded asset $A_0$ and
liability $L_0$ completely across the same item population. Those allocations
must sum exactly to the book's recorded gross amounts. They explain inherited
balances rather than inventing a new assessment of the original opening.
For the supported profit-or-loss route, the change in deferred expense is
$(L_1-L_0)-(A_1-A_0)$. The journal debits the asset by $A_1-A_0$, credits the
liability by $L_1-L_0$, and records the balancing expense or benefit. Negative
changes reverse the corresponding balance. The request must support that
allocation; an equity, other-comprehensive-income or business-combination route
is refused by this particular implementation instead of being charged to profit.

In the checked example, property, plant and equipment has carrying amount 100
and tax base 60. At rate $1/4$, its taxable difference of 40 gives a liability
of 10. Debt has carrying amount 100 and tax base 60, so its difference is
deductible. Recognizing 32 of the available 40 gives an asset of 8 and leaves 8
unrecognized. Relative to the recorded asset of 10 and liability of 4, the asset
falls by 2 and the liability rises by 6; deferred expense is therefore 8.
Cash stays $721/2$, current tax and its unpaid liability are unchanged, and
recorded profit becomes $-33/4$. Assets of $1397/2$ equal liabilities of $459/4$
plus equity of $2335/4$. Adding the deferred expense back in the indirect cash
calculation leaves operating cash at $-59/2$.

Loss recognition uses the same capacity discipline. Recognizing deductions of
32 for the debt and 10 for a loss lot needs capacity of at least 42; capacity
of 32 is insufficient even if each item separately fits it. With capacity 42
and rate $1/4$, their combined deferred asset is $21/2$. The loss ledger still
contains the original remaining quantity of 19: recognizing a tax benefit does
not use the loss in the current-tax calculation. Matching bases, zero rates,
partial recognition and opposite asset/liability differences have separate
checked cases.

The proposal retains the entire prior result and changes only deferred balances,
deferred expense and its new event identity. Cash flows, current tax, loss
history, component state and assessment dates remain exact. Unknown necessary
findings stay unresolved, and an inconsistent population, allocation, base or
capacity produces no changed financial state. Pure apply compares the complete
recomputed proposal. Actual source and accounting admission, broader recovery
and allocation methods and annual completeness remain open. The preceding
frozen carry interface accepts a current-tax reassessment request; consuming
this newer result requires an explicit successor interface that preserves the
deferred journal and recognition schedule.
